(1) 一个电荷振子质量 \(m>0\)、固有角频率 \(\omega_0>0\)，在非相对论经典Abraham–Lorentz近似下满足
\[
m\ddot x+m\omega_0^2x=m\tau\dddot x,\qquad \tau>0,\quad \tau\omega_0\ll1.
\]
建立指数解 \(x=e^{\lambda t}\) 的特征三次方程，证明存在一个正实根（跑逸根），求其 \(\tau\to0\) 的主导标度，说明其不应作为物理可接受的自由初值响应。
(2) 通过对小 \(\tau\) 的降阶，把 \(\dddot x\) 在力项中用零阶振子方程代替，写出稳定二阶有效方程及复频率的首阶修正，求振幅与机械能的长期衰减率。
(3) 不采用有效黏性阻尼猜测，直接以原三阶方程乘 \(\dot x\)，推导机械能与可交换的近场Schott项、向外辐射的Larmor功率之间的严格能量恒等式，说明Schott项符号与辐射功率始终非负。
(4) 用降阶近似的自由谐振运动求周期平均辐射功率，独立检验机械能衰减率；讨论 \(\tau\to0\) 极限、经典点电荷近似、预加速与外驱动因果处理的适用边界。
